\section{Introduction}
A fixed per-image channel-symbol budget must accommodate source information, signaling, and protection. Digital coding and learned joint source--channel coding (JSCC) address this constraint with different representations~\cite{deepjscc,swinjscc}. At a generative receiver, transmitted information constrains reconstruction to the source while the prior predicts missing content.

Next-scale tokenization provides complete coarse maps covering the latent grid, followed by residual refinements. Visual autoregressive modeling (VAR) reported state-of-the-art results in its original class-conditional ImageNet generation benchmarks~\cite{var}, motivating this completion prior. ARPC uses next-scale entropy coding and completion~\cite{arpc}; Ada-TokenCom combines next-token prefix coding, tail generation, and wireless adaptation~\cite{adatokencom}. TONIC considers task-oriented token communication~\cite{tonic}. Here, the task is source-image reconstruction under a hard per-image symbol budget.

Fitting progressive descriptions to channel resources is classical~\cite{chande}. A complete next-scale residual map can exceed remaining frame capacity even when some tokens fit. Image-dependent entropy lengths require actual decodable length checks, and incomplete boundary maps must enter both residual reconstruction and subsequent prediction.

NeST-Com combines token-wise endpoints, independently terminated arithmetic streams, and prefix-constrained reconstruction. Shared coarse maps determine coding probabilities; received tokens overwrite predicted positions~\cite{var}, updating the residual latent and finer-scale context. Shared-model controls favor NeST over whole-scale transmission on DISTS and DreamSim at all seven budget/SNR points. Matched grids and receiver interventions examine the design, while external comparisons establish complete-system tradeoffs.
